\section{Time Differentiation of Operators}

\SourcePage{42}{45}
In quantum mechanics, the concept of the time derivative of a physical
quantity cannot be defined with the sense that it carries in classical
mechanics. Indeed, the
definition of the derivative in classical mechanics is connected with
considering the values of the quantity at two close but different instants.
But in quantum mechanics a quantity that has a definite value at some instant
has no definite value at all at subsequent instants; this was discussed
in more detail in \secref{1}.

For this reason the concept of the time derivative has to be given a
different definition in quantum mechanics. A natural course is to define the
derivative $\dot f$ of a quantity $f$ as that quantity whose mean value
equals the time derivative of the mean value $\bar f$. Thus, by definition,
we have
\begin{equation}
  \bar{\dot f}=\dot{\bar f}.
  \tag{9.1}
\end{equation}
Taking this definition as the starting point, it is not hard to derive an
expression for the quantum-mechanical operator $\hat{\dot f}$ that
corresponds to%
\SourcePage{43}{46}
the quantity $\dot f$:
\[
  \begin{aligned}
    \bar{\dot f}=\dot{\bar f}
      &=\frac{d}{dt}\int\Psi^*\hat f\Psi\,dq={}\\
      &=\int\Psi^*\frac{\partial\hat f}{\partial t}\Psi\,dq
       +\int\frac{\partial\Psi^*}{\partial t}\hat f\Psi\,dq
       +\int\Psi^*\hat f\frac{\partial\Psi}{\partial t}\,dq.
  \end{aligned}
\]
In this expression $\partial\hat f/\partial t$ stands for the operator that
results from differentiating the operator $\hat f$ with respect to time, on
which the latter may depend parametrically. Substituting for the derivatives $\partial\Psi/\partial t$,
$\partial\Psi^*/\partial t$ their expressions according to (8.1), we obtain
\[
  \bar{\dot f}
  =\int\Psi^*\frac{\partial\hat f}{\partial t}\Psi\,dq
  +\frac{i}{\hbar}\int(\hat H^*\Psi^*)\hat f\Psi\,dq
  -\frac{i}{\hbar}\int\Psi^*\hat f(\hat H\Psi)\,dq.
\]
Since the operator $\hat H$ is Hermitian,
\[
  \int(\hat H^*\Psi^*)(\hat f\Psi)\,dq
  =\int\Psi^*\hat H\hat f\Psi\,dq;
\]
thus, we have
\[
  \bar{\dot f}
  =\int\Psi^*\left(
    \frac{\partial\hat f}{\partial t}
    +\frac{i}{\hbar}\hat H\hat f
    -\frac{i}{\hbar}\hat f\hat H
  \right)\Psi\,dq.
\]

Since, on the other hand, by the definition of mean values we must have
$\bar{\dot f}=\int\Psi^*\hat{\dot f}\Psi\,dq$, it is seen that the
expression in parentheses under the integral is the required operator
$\hat{\dot f}$\footnote{In classical mechanics, for the total time
derivative of a quantity $f$ that is a function of the generalized coordinates
$q_i$ and momenta $p_i$ of the system, we have
\[
  \frac{df}{dt}=\frac{\partial f}{\partial t}
  +\sum_i\left(
    \frac{\partial f}{\partial q_i}\dot q_i
    +\frac{\partial f}{\partial p_i}\dot p_i
  \right).
\]
Substituting, according to Hamilton's equations,
$\dot q_i=\partial H/\partial p_i$, $\dot p_i=-\partial H/\partial q_i$, we
obtain
\[
  \frac{df}{dt}=\frac{\partial f}{\partial t}+[H,f],\qquad
  [H,f]\equiv\sum_i\left(
    \frac{\partial f}{\partial q_i}\frac{\partial H}{\partial p_i}
    -\frac{\partial f}{\partial p_i}\frac{\partial H}{\partial q_i}
  \right);
\]
$[H,f]$ is the so-called \textit{Poisson bracket} for the quantities $f$ and
$H$ (see I, \S~42). Comparing with expression (9.2), we see that in the
transition to the classical limit the operator
$i(\hat H\hat f-\hat f\hat H)$ becomes zero to the first approximation, as it
should, and in the next approximation (in $\hbar$) becomes the quantity
$\hbar[H,f]$.

This result is also valid for any two quantities $f$ and $g$: in the limit the
operator $i(\hat f\hat g-\hat g\hat f)$ becomes the quantity $\hbar[f,g]$,
where $[f,g]$ is the Poisson bracket
\[
  [f,g]\equiv\sum_i\left(
    \frac{\partial g}{\partial q_i}\frac{\partial f}{\partial p_i}
    -\frac{\partial g}{\partial p_i}\frac{\partial f}{\partial q_i}
  \right).
\]
This is a consequence of the fact that we can always, formally, conceive of a
system whose Hamiltonian coincides with $\hat g$.}:
\begin{equation}
  \hat{\dot f}
  =\frac{\partial\hat f}{\partial t}
  +\frac{i}{\hbar}(\hat H\hat f-\hat f\hat H).
  \tag{9.2}
\end{equation}
\SourcePage{44}{47}
If the operator $\hat f$ does not depend explicitly on time, then, up to a
factor, $\hat{\dot f}$ reduces to the commutator of the operator $\hat f$
with the Hamiltonian.

The physical quantities whose operators have no explicit dependence on time
and which, moreover, commute with the Hamiltonian---so that
$\hat{\dot f}=0$---form a very important category. Such quantities are called
\textit{conserved}. For these,
$\bar{\dot f}=\dot{\bar f}=0$, i.e.\ $\bar f=\mathrm{const}$. Put
differently, the quantity's mean value stays constant in time. One can
assert, moreover, that if the quantity $f$ has a definite value in a given
state (i.e.\ if the wave function is an eigenfunction of the operator
$\hat f$), then at later instants as well it will have a definite value---and
the very same one.
